
\subsubsection{4.1 Research Questions}

\textbf{RQ1:} Does step-local grounded DPO feedback (Condition A) produce higher locality scores than permuted control (Condition P), confirming the oracle mechanism at the probability mass level? [Tests H-E1]

\textbf{RQ2:} Does the DPO pair construction protocol achieve 100\% state alignment? [Tests H-M1]

\textbf{RQ3:} Does the oracle locality effect translate to greater hard-stratum pass@1 recovery, and does this advantage vary monotonically with α? [Tests H-M3, H-M4; not executed]

\subsubsection{4.2 Datasets}

\textbf{miniF2F Hard Subset.} Full benchmark: 488 problems (AMC/AIME/IMO and high-school/undergraduate mathematics) in Lean4 format, loaded via HuggingFace (\texttt{Tonic/MiniF2F}). Format: \texttt{{name, split, formal_statement, goal, header}}. The hard subset is defined as problems where BFS-Prover cold-start SFT achieves pass@1 < 20\% across 16 rollouts. The full miniF2F dataset (488 problems) was loaded from HuggingFace in the Phase 4 run; the hard subset construction required cold-start SFT evaluation that was not completed (see Section 5.2).

\textbf{Vericoding Hard Subset.} Full benchmark: 12,504 formally verified code specifications across Lean4/Verus/Dafny \citep{Bursucetal2025}. The Vericoding data path was not found during the Phase 4 run (\texttt{data/vericoding} absent); 0 Vericoding problems were loaded. Any Vericoding results reported below are artifacts of the same synthetic fallback that produced the miniF2F null result.

\begin{table}[htbp]
\centering
\begin{tabular}{llll}
\toprule
Dataset & Full Size & Intended Hard Subset & Actual Status in Phase 4 \\
\midrule
miniF2F & 488 problems & ~100–150 (pass@1 < 20\%) & 488 loaded; hard subset not constructed \\
Vericoding (Lean4) & ~3,000 (estimated) & ~300–600 & 0 loaded; data path not found \\
\bottomrule
\end{tabular}
\end{table}

\subsubsection{4.3 Model}

\textbf{Base model:} \texttt{ByteDance-Seed/BFS-Prover-V2-7B} — Qwen2.5-Math-7B with BFS-Prover cold-start SFT initialization. Architecture: decoder-only transformer, 7B parameters, bfloat16. The model serves as both the frozen reference π_ref and the starting point for each condition's DPO fine-tuning. The model was loaded with \texttt{device_map="auto"} on a single GPU (GPU 2 per the experiment log), using 90\% of device memory.

\subsubsection{4.4 DPO Training Configuration}

\begin{table}[htbp]
\centering
\begin{tabular}{lll}
\toprule
Hyperparameter & Value & Source \\
\midrule
DPO β & 10.0 & BFS-Prover [Xin et al., 2025] \\
Learning rate start & 5e-6 & BFS-Prover \\
Learning rate end & 5e-7 (linear decay) & BFS-Prover \\
Batch size & 16 & BFS-Prover \\
Epochs & 1 & Standard DPO practice \\
Optimizer & AdamW (wd=0.01, β=(0.9, 0.999)) & eric-mitchell/DPO \\
Mixed precision & bfloat16 & BFS-Prover \\
Seeds & 1 & PoC stage \\
\bottomrule
\end{tabular}
\end{table}

The DPO loss uses \texttt{average_log_prob=False} (sum log-probabilities for variable-length tactic sequences), consistent with the eric-mitchell/DPO reference implementation. The concatenated forward pass pattern (chosen and rejected in a single forward pass) was implemented as specified.

\subsubsection{4.5 Infrastructure}

Single NVIDIA H100 GPU (80GB), CUDA 12.0+, Python 3.10+. LeanDojo (\texttt{lean-dojo>=2.0.0}) was listed as a required dependency but was not available at runtime due to the absent Lean4 toolchain on the H100 cluster.

\subsubsection{4.6 Evaluation Metrics}

\textbf{Primary metric (H-E1 gate):} Locality Score (LS) — tested via one-sided t-test (\texttt{scipy.stats.ttest_1samp}, alternative='greater'), significance threshold p < 0.05. Gate: LS_A > LS_P.

\textbf{Secondary metric:} LS_A > LS_B (grounded versus ungrounded comparison).
